\section{Recomendaciones prácticas y operación del sistema}
\subsection{Tablero de monitoreo de Autonomy}
El ponente recomienda construir un tablero con tres indicadores, reward accuracy, reset success y goal coverage, para la revisión diaria.
\begin{tabularx}{\textwidth}{lX}
\toprule
Dimensión & Indicador \\
\midrule
Reward fidelity & classifier precision/recall, clip similarity drift; \\
Reset coverage & tasa de éxito reciente de forward-backward > 80\%; \\
Goal expansion & número de goals activos y goals nuevos por semana; \\
\bottomrule
\end{tabularx}

\subsection{Ciclo de retroalimentación entre etapas}
En cada ejecución de un experimento de Autonomy se registra lo siguiente:
\begin{enumerate}
\item la versión actual del modelo de reward y de la safety penalty;
\item el failure pattern de la Reset policy;
\item la incorporación de nuevos objetivos o habilidades y las decisiones de retirar objetivos anteriores;
\end{enumerate}
Esta información forma un cross-team log que facilita la replicación futura.

\begin{knowledgebox}{Plantilla de Feedback Loop}
1. Pre-scale: confirmar el estado de reward, safety y checkpoint;\\
2. Scale: ejecutar la policy y sincronizar las metrics;\\
3. Post-scale: analizar en retrospectiva el metrics drift y, si se requiere un rollback, activarlo de inmediato.
\end{knowledgebox}

\subsection{Resumen de la sección}
En la práctica, Autonomy requiere un tablero, retroalimentación entre etapas y un plan de reversión para que los resultados de la investigación puedan desplegarse una y otra vez.
\newpage
